
\section{Зорилго}
	\quad \quad	МУИС-ийн Дипломын ажлыг удирдах системийн хэрэглэгчийн интерфэйс дээр React болон JSF технологийг хослуулсан микро-фронтэнд архитектурыг турших.

\section{Зорилт}
	\quad \quad	Дээрх зорилгод хүрэхийн тулд дараах зорилтуудыг дэвшүүлэн ажиллана. Үүнд:
	\begin{itemize}
		\item Микро-фронтэнд архитектур болон React, JSF технологиудыг нэгтгэх интеграцын аргуудын талаарх өмнөх судалгаануудыг харьцуулан судлах;
		\item React, JSF технологийг уялдуулан ажиллуулах микро-фронтэнд системийн архитектурыг зохиомжлох;
		\item Загварчилсан архитектурын дагуу системийн бодит туршилтын загварыг хөгжүүлэх;
		\item Хөгжүүлсэн холимог системдээ программ хангамжийн чанарын үзүүлэлтүүдийн үнэлгээ хийж, уламжлалт арга барилтай харьцуулсан дүгнэлт гаргана.
	\end{itemize}

\section{Сэдэв сонгосон үндэслэл}
	\quad \quad	Шинээр системүүд хөгжүүлэхэд хэрэглэгчийн шаардлагад нийцүүлэн орчин үеийн технологи ашиглах хэрэгцээ гардаг хэдий ч, эдгээрийг одоогийн дэд бүтэцтэй нийцтэйгээр уялдуулан хөгжүүлэх шаардлага тулгардаг.Үүнийг шийдвэрлэх үр дүнтэй арга бол Микрофронтэнд архитектур юм. Энэхүү архитектур нь өөр өөр технологиор хөгжүүлэгдсэн хэсгүүдийг нэгтгэн, нэг цогц систем болгон ажиллуулах боломжийг олгодог. Тиймээс энэхүү судалгааны ажлын хүрээнд МУИС-ийн "Дипломын ажлыг удирдах систем"-ийн туршилтын загварыг шинээр зохиомжлохдоо React, JSF технологийг хослуулсан микро-фронтэнд архитектурыг ашиглахаар зорьсон болно. Энэ нь орчин үеийн болон уламжлалт технологиудыг хэрхэн амжилттай уялдуулан ажиллуулж болохыг судлах, түүнчлэн энэхүү холимог шийдэл нь программ хангамжийн чанарын үзүүлэлтүүдэд хэрхэн нөлөөлж буйг бодитоор үнэлж дүгнэх чухал ач холбогдолтой юм.
\section{Сэдвийн судлагдсан байдал}

\section{Судалгааны ажлын шинэлэг тал}

\section{Судалгааны ажлын үр дүн, ач холбогдол}
